% ============================================
% MEM-04
% Gerado de molambudos.md — build_miolo.py
% ============================================
\fragtitle{mem}{M}{MEM}{04}{A Vala}
\markboth{MEM-04 --- A Vala}{}

\noindent \textit{O encontro}

\noindent \textbf{1915 --- Na rota de Senador Pompeu a Fortaleza}

\vspace{0.5em}\noindent\rule{\textwidth}{0.5pt}\vspace{0.5em}

\noindent Meu pai morreu de manhã. O sol nasceu e ele já estava duro, com os olhos abertos para o céu que não tinha nuvem. Minha mãe passou a mão nos olhos dele e disse: \textit{\textquotedblleft{}Anda, Joaquim. Teu pai já foi.\textquotedblright{}}

\noindent Eu tinha oito anos.

\noindent A gente vinha andando há quatro dias. Antes éramos seis. Agora éramos três: eu, minha mãe e o velho Anselmo, que não era parente mas caminhava junto porque sozinho ninguém sobrevive. O velho Anselmo disse que tinha visto a lagoa secar em 1877. Disse que aquilo sim foi seca. Eu olhei para ele e vi que ele não acreditava no que dizia. Os olhos dele estavam iguais aos do meu pai.

\noindent No quinto dia, a gente encontrou o comboio.

\noindent Eram soldados. Magros, sujos, com fardas remendadas. Um deles levava uma espingarda. O outro levava um chicote.

\noindent \textit{\textquotedblleft{}Pra onde vão?\textquotedblright{}}, perguntou o do chicote.

\noindent Minha mãe respondeu que íamos para o litoral. Que ouviu dizer que lá tinha comida.

\noindent O soldado riu. \textit{\textquotedblleft{}Comida? O litoral tá cheio de retirante que nem você. Tá todo mundo morrendo na praia.\textquotedblright{}}

\noindent Minha mãe não respondeu.

\noindent \textit{\textquotedblleft{}Vem com a gente\textquotedblright{}\textit{, disse o soldado. }\textquotedblleft{}Tem o Campo do Alagadiço, perto de Fortaleza. Lá dá comida.\textquotedblright{}}

\noindent Ela hesitou. Eu senti a mão dela apertar a minha.

\noindent \textit{\textquotedblleft{}É ordem\textquotedblright{}}, disse o soldado, e o chicote estalou no ar. Não bateu em ninguém. Mas o som ficou. Como se tivesse marcado a gente.

\vspace{0.5em}\noindent\rule{\textwidth}{0.5pt}\vspace{0.5em}

\noindent O Curral do Governo era uma cerca de arame farpado em volta de um descampado. Não tinha árvore. Não tinha sombra. Tinha gente. Muita gente. Mais gente do que eu tinha visto em toda a minha vida.

\noindent Eles estavam deitados no chão, encostados uns nos outros. Alguns mexiam a boca sem fazer som. Outros não mexiam nada.

\noindent \textit{\textquotedblleft{}Molambudos\textquotedblright{}}, alguém disse.

\noindent Eu não sabia o que significava. Depois eu soube. Significa farrapos. Trapos humanos. Gente que não é mais gente porque a seca levou tudo o que tinha dentro.

\noindent Ficamos ali. Minha mãe achou um canto perto da cerca. O velho Anselmo sentou e não levantou mais. No segundo dia, ele disse: \textit{\textquotedblleft{}Tô indo.\textquotedblright{}} E foi.

\noindent Minha mãe começou a tossir.

\vspace{0.5em}\noindent\rule{\textwidth}{0.5pt}\vspace{0.5em}

\noindent A vala ficava atrás do curral, a uns duzentos metros. Toda noite os soldados carregavam os corpos para lá. Durante o dia dava pra ver os urubus. Eles faziam um círculo no céu, rodando, rodando, como se o céu fosse um olho gigante.

\noindent Eu fui até a vala numa noite em que minha mãe tossia sem parar.

\noindent Não fui por curiosidade.
Fui porque ouvi uma coisa.

\noindent Era uma voz.

\noindent Não uma voz humana --- mas uma voz. Ela cantava. Não uma canção com letra. Uma melodia. Baixa. Rouca. Uma melodia que subia e descia como o vento num lugar vazio.

\noindent Ela cantava:

\noindent \textit{\textquotedblleft{}Quem tem fome senta. Quem tem fome espera. Quem tem fome come o que vier.\textquotedblright{}}

\noindent Eu parei na borda da vala.

\noindent Lá embaixo, entre os corpos, algo se mexeu.

\noindent Eu devia ter corrido. Devia ter voltado para o curral. Devia ter enterrado o rosto na saia da minha mãe e esquecido que tinha ouvido alguma coisa.

\noindent Mas eu era uma criança.
E crianças que passam fome não fogem de comida.
Elas investigam.

\vspace{0.5em}\noindent\rule{\textwidth}{0.5pt}\vspace{0.5em}

\noindent A vala não era funda. Uns três metros. Os corpos se empilhavam como sacos de farinha. Alguns estavam nus. Outros amarrados pelas mãos. Os olhos deles estavam abertos. Todos os olhos estavam abertos.

\noindent No meio da pilha, algo se moveu com cuidado.

\noindent Não era o movimento de quem luta para sair.
Era o movimento de quem está \textit{arrumando} os corpos à sua volta.

\noindent Um braço emergiu. Não um braço humano --- era fino demais, longo demais, com articulações onde não devia ter. A mão tinha dedos demais, e cada dedo tinha uma unha preta que brilhava como vidro molhado.

\noindent A mão pegou um braço solto na pilha e o colocou de lado. Com cuidado. Como quem arruma a mesa antes de uma refeição.

\noindent Outro braço surgiu. Depois a cabeça.

\noindent A cabeça era quase humana. Quase. A forma era de uma caveira, mas a pele ainda cobria o osso --- uma pele fina, cinzenta, esticada como papel molhado sobre os ossos do rosto.

\noindent O nariz era um buraco.
A boca era uma linha.
Os olhos eram duas fendas.

\noindent E dentro das fendas, um brilho amarelo.

\noindent Ela me olhou.

\noindent Não se mexeu. Não fez nada. Apenas me olhou, com a cabeça inclinada, como um cachorro ouvindo um som estranho.

\noindent Eu não corri.

\noindent Ela começou a subir.

\noindent Não como uma pessoa sobe. Como uma aranha. Os braços e as pernas se moviam em coordenação, encontrando apoio nos corpos, deslizando entre eles sem pressa. Os ossos estalavam sob o peso dela. A carne cedia.

\noindent Ela saiu da vala.

\noindent Ficou de pé na minha frente.

\noindent Ela era imensa --- não alta, mas \textit{larga}. O corpo era uma gaiola de costelas viradas para fora, como se o esqueleto tivesse crescido pelo avesso. Entre as costelas, algo pulsava. Um órgão grande, escuro, coberto de veias que brilhavam como minhocas.

\noindent O cheiro dela chegou até mim.

\noindent O cheiro não vinha da morte. Vinha de \textit{comida}. Carne assada. Caldo grosso. O cheiro que saía da cozinha de Dona Sinhá antes da seca. O cheiro que eu não sentia há meses.

\noindent Meu estômago doeu.

\noindent \textit{\textquotedblleft{}Cê tem fome?\textquotedblright{}}

\noindent A voz dela era macia. Quase doce. Não combinava com o corpo. Era a voz de uma avó perguntando se o neto já comeu.

\noindent Eu não respondi.

\noindent \textit{\textquotedblleft{}Cê tem.\textquotedblright{}\textit{ Ela inclinou a cabeça para o outro lado. }\textquotedblleft{}Eu sinto. Tua fome é grande. Grande que nem a minha.\textquotedblright{}}

\noindent \textit{\textquotedblleft{}A senhora...\textquotedblright{}\textit{ Minha voz saiu fraca. }\textquotedblleft{}A senhora é o quê?\textquotedblright{}}

\noindent Ela pensou. O órgão entre as costelas pulsou mais forte.

\noindent \textit{\textquotedblleft{}Sou o que sobra\textquotedblright{}\textit{, ela disse. }\textquotedblleft{}Sou o que cresce onde muito morre. Sou o que a terra não consegue engolir.\textquotedblright{}}

\noindent \textit{\textquotedblleft{}A senhora come os mortos?\textquotedblright{}}

\noindent \textit{\textquotedblleft{}Alguém tem que comer.\textquotedblright{}}

\noindent Ela disse isso com naturalidade. Como quem diz \textquotedblleft{}o sol nasce todo dia\textquotedblright{}.

\noindent \textit{\textquotedblleft{}Os vivos não querem. Os vivos enterram. Os vivos choram. Depois vão embora. Os mortos ficam. Alguém tem que aproveitar.\textquotedblright{}}

\noindent Ela estendeu a mão para mim. A mão grande, de dedos longos, com unhas pretas.

\noindent \textit{\textquotedblleft{}Cê quer?\textquotedblright{}}

\noindent \textit{\textquotedblleft{}Quer o quê?\textquotedblright{}}

\noindent \textit{\textquotedblleft{}Comer.\textquotedblright{}}

\noindent Eu olhei para a mão dela. Olhei para a vala. Olhei para os corpos.

\noindent \textit{\textquotedblleft{}Não\textquotedblright{}\textit{, eu disse. }\textquotedblleft{}Não quero comer eles.\textquotedblright{}}

\noindent Ela riu. A risada parecia o som de folhas secas se mexendo.

\noindent \textit{\textquotedblleft{}Não é eles que eu tô oferecendo, menino.\textquotedblright{}}

\noindent \textit{\textquotedblleft{}Então o quê?\textquotedblright{}}

\noindent Ela se abaixou. O rosto ficou na altura do meu. O cheiro de comida ficou mais forte. Meu estômago doeu de novo.

\noindent \textit{\textquotedblleft{}Eu como por você\textquotedblright{}\textit{, ela disse. }\textquotedblleft{}Eu como todos os dias. Eu como os que morrem. Eu como os que ninguém quis. E você nunca mais vai sentir fome.\textquotedblright{}}

\noindent Ela esperou.

\noindent \textit{\textquotedblleft{}Só que...\textquotedblright{}}

\noindent \textit{\textquotedblleft{}Só que o quê?\textquotedblright{}}

\noindent Ela sorriu. A boca se abriu mais do que devia. Os dentes estavam em fileiras demais. Fileiras que não cabiam numa boca humana.

\noindent \textit{\textquotedblleft{}A fome não vai embora. Ela fica. Só que ela não é mais tua. É minha. E eu vou morar dentro de ti para sentir ela.\textquotedblright{}}

\noindent Ela tocou meu peito com um dedo.

\noindent \textit{\textquotedblleft{}Tu aceita?\textquotedblright{}}

\noindent Eu devia ter dito não.

\noindent Eu devia ter cuspido na mão dela. Devia ter corrido. Devia ter gritado.

\noindent Mas eu tinha oito anos.
E tinha tanta fome.
Tanta fome que o não não saía.

\noindent \textit{\textquotedblleft{}Sim\textquotedblright{}}, eu disse.

\noindent Ela abriu a boca.

\noindent Não foi um beijo. Não foi uma mordida. Foi um sopro. Um sopro quente que saiu da boca dela e entrou na minha. Desceu pela garganta. Encheu o peito. Escorreu para o estômago.

\noindent Ela estava dentro de mim.

\noindent \textit{\textquotedblleft{}Pronto\textquotedblright{}\textit{, ela disse. }\textquotedblleft{}Agora a fome é minha.\textquotedblright{}}

\noindent Ela se levantou. Olhou para o céu. A lua brilhava.

\noindent \textit{\textquotedblleft{}Vai\textquotedblright{}\textit{, ela disse. } \textquotedblleft{}Vai viver. Depois a gente se encontra de novo.\textquotedblright{}}

\noindent Ela voltou para a vala. Desceu entre os corpos. Deitou. E fechou os olhos.

\noindent Eu fiquei parado na borda, olhando.

\noindent Não sentia mais fome.
Mas sentia um peso.
Um peso no fundo do estômago.
Um peso que respirava.\NE{O trecho do pacto com a vala é o único em todo o diário em que Joaquim escreve com a mão esquerda. Os peritos que examinaram o manuscrito notaram: a caligrafia muda abruptamente na linha \textquotedblleft{}Sim, eu disse\textquotedblright{}, a pressão da caneta sobre o papel é diferente, e o traçado das letras se desvia para a esquerda — como se outra mão estivesse guiando a dele. O editor não interveio. A escolha de manter a caligrafia original é uma escolha ética e documental.}

\vspace{0.5em}\noindent\rule{\textwidth}{0.5pt}\vspace{0.5em}

\noindent Naquela noite, na beira de uma vala de mortos, no meio de um campo de concentração, um menino de oito anos fez um trato com um monstro.

\noindent O monstro comeu a fome dele.

\noindent Mas deixou o buraco.

\noindent O buraco ficou.

\noindent E tudo o que entrou no buraco depois --- anos, sangue, gente, memórias --- nunca encheu.

\noindent Porque o buraco não é meu.

\noindent É dela.

\noindent Ela mora dentro de mim desde aquela noite.

\noindent Ela come. Eu sinto.

\noindent Ela dorme. Eu acordo.

\noindent Ela espera. Eu vivo.

\noindent Ela está esperando você também.

\vspace{0.5em}
\noindent\textit{\textbf{↪ Rotas:} \rota{MEM-05} \rotasep \rota{DOC-01} \rotasep \rota{LUC-01}}